\section{Elementary inputs for the actual root-block mass theorem}
\label{ebi-section}

For the actual blocks in Section~\ref{fml-section}, the archimedean
and two fixed-prime inputs have an elementary replacement when $B\ge n$.
This gives a proof of their full-mass conclusion without logarithmic
forms. It does not estimate the signed cost of the remaining primes.
Throughout this section $B\ge n\ge1$ and $B\le k<2B$ are integers,
\[
 a=3k,\quad w=a+\zeta,\quad Q=a^2+a+1,\quad
 T_d=|P(w^d)|,\quad t=\log c_n,\quad \Delta=3t-\log T_n.
\]
The auxiliary index $d$ is a positive integer. Here $c_n$ is the largest absolute boundary arm. The elementary norm
bounds give
\begin{equation}\label{ebi-norm}
 Q^{n/2}\le c_n\le\frac2{\sqrt3}Q^{n/2},\qquad
 t\ge n\log(3B),\qquad \log T_d\le3d\log(6B+1).
\end{equation}
These primitive unramified nonunit roots have nonzero boundaries;
for the index $n$ this also follows directly from the angle below.

\begin{lemma}[A direct angular estimate]\label{ebi-angle}
On every such actual block,
\begin{equation}\label{ebi-angular-bound}
 T_n\ge\frac{2n}{\pi(12B+1)}Q^{3n/2},\qquad
 0\le\frac\Delta t\le\frac4n.
\end{equation}
\end{lemma}
\begin{proof}
Write $w=\sqrt Q\,e^{i\theta}$, where
$\theta=\arctan(\sqrt3/(6k+1))$. The elementary bounds
$x/2\le\arctan x\le x$ for $0\le x\le1$ imply
\[
 \frac{\sqrt3}{2(12B+1)}\le\theta,\qquad
 0<3n\theta<\frac{\sqrt3}{2}<\frac\pi2.
\]
Thus $\sin(3n\theta)\ge3n\sqrt3/(\pi(12B+1))$.
The exact identity
$z^3-\bar z^3=3\sqrt{-3}\,P(z)$ gives
\[
 T_n=\frac2{3\sqrt3}Q^{3n/2}\sin(3n\theta),
\]
proving the first bound. Every arm has absolute value at most $c_n$,
so $\Delta\ge0$. Combining the first bound with
\eqref{ebi-norm} gives
\[
 \Delta\le\log\frac{4\pi(12B+1)}{3\sqrt3\,n}
          \le4\log(3B).
\]
For the last inequality, $12B+1\le13B$, $\pi<4$ and $\sqrt3>1$
bound the logarithm's argument by $81B/n\le(3B)^4$.
Division by $t\ge n\log(3B)$ completes the proof.
\end{proof}

\begin{lemma}[Exact fixed-prime depths]\label{ebi-fixed-depth}
For all the above roots and indices,
\begin{equation}\label{ebi-three}
 v_3(T_n)=v_3(T_1)+v_3(n),
\end{equation}
\begin{equation}\label{ebi-two}
 v_2(T_n)=
 \begin{cases}
  v_2(T_1),&n\text{ odd},\\
  v_2(T_2)+v_2(n/2),&n\text{ even}.
 \end{cases}
\end{equation}
Consequently
\begin{equation}\label{ebi-small-cost}
 \frac{v_2(T_n)\log2+v_3(T_n)\log3}{t}\le\frac{20}{n}.
\end{equation}
\end{lemma}
\begin{proof}
We first supply the lifting step. Let $v$ be a nonarchimedean
valuation above $p$, normalized by $v(p)=1$, and write $u=1+\delta$.
If $e=v(\delta)>1/(p-1)$, the linear term in
$(1+\delta)^p-1$ has valuation $1+e$. Every intermediate term
has valuation $1+je>1+e$, and the last has valuation $pe>1+e$.
Thus $v(u^p-1)=e+1$. For an exponent $r$ prime to $p$, the weaker
condition $e>0$ suffices: the linear term has valuation $e$ and all
higher terms have larger valuation. Hence $v(u^r-1)=e$.
Iteration proves $v(u^m-1)=v(u-1)+v_p(m)$ under the first threshold.
This binomial argument applies also at the ramified quadratic place.

Put $\eta=(w/\bar w)^3$. Since $Q$ is prime to both two and three,
$w,\bar w$ are units there, and
\begin{equation}\label{ebi-local-identity}
 \eta^d-1=\frac{3\sqrt{-3}\,P(w^d)}{\bar w^{3d}}.
\end{equation}
At three its valuation is $v_3(T_d)+3/2$. The integer
$T_1=a(a+1)$ is divisible by three, so $v_3(\eta-1)\ge5/2>1/2$.
Apply the lifting step and subtract $3/2$ to obtain
\eqref{ebi-three}.

At two the constant $3\sqrt{-3}$ is a unit and $T_1$ is even,
so the odd-exponent assertion follows from the weaker lifting step.
For an even exponent, direct multiplication gives
\begin{equation}\label{ebi-t2}
 T_2=a(a-1)(a+1)(a+2)(2a+1).
\end{equation}
The four consecutive factors $a-1,a,a+1,a+2$ have product divisible
by eight. Thus $v_2(\eta^2-1)=v_2(T_2)\ge3>1$.
Apply the stronger lifting step to $\eta^2$ at exponent $n/2$ to
obtain \eqref{ebi-two}.

Write $L=\log(6B+1)$. Equations~\eqref{ebi-norm}--\eqref{ebi-two}
bound the two-prime cost by $9L+2\log n$; the odd case even has
a smaller bound. Since $L\le2\log(3B)$ and
$\log n\le\log(3B)$, this is at most $20\log(3B)$.
Normalize by $t\ge n\log(3B)$.
\end{proof}

\begin{theorem}[An elementary replacement in the mass-location chain]
\label{ebi-location}
Set $B=n^4$,
\[
 Z=\left\lfloor\frac{B\sqrt{\log n}}{1+\log\log n}\right\rfloor,
 \qquad \eta_n=(\log n)^{-1/4}.
\]
For all sufficiently large integers $n$, outside at most a proportion
$C\eta_n$ of the actual roots in the block,
\begin{equation}\label{ebi-tail}
 3-\frac4n-\eta_n
 \le\frac1{t_k}\sum_{p>Z}v_p(T_n(k))\log p\le3.
\end{equation}
If $n$ is prime, the same roots also satisfy
\begin{equation}\label{ebi-top}
 3-\frac7n-\eta_n
 \le\frac1{t_k}\sum_{\substack{p>Z\\d_p(k)=n}}
       v_p(T_n(k))\log p\le3.
\end{equation}
Every included valuation in \eqref{ebi-top} is its top-rank first depth.
\end{theorem}
\begin{proof}
The full-mass interval estimate of Section~\ref{fml-section} for primes
at least five uses actual rank counts, a harmonic sum and
Brun--Titchmarsh. Its all-index specialization has mean
$O((\log n)^{-1/2})$, by the proved elementary divisor estimates.
Add \eqref{ebi-small-cost}. Markov at $\eta_n$ bounds the full
small-prime mass by $\eta_n t_k$ outside the stated proportion.
Subtract this mass from $\log T_n=3t_k-\Delta$ and apply
\eqref{ebi-angular-bound}, obtaining \eqref{ebi-tail}.
At prime index the only ranks are one and $n$. Above $Z>n$ there
is no exponent lifting, and the entire rank-one mass is at most
$\log T_1\le3t_k/n$. Subtract it to obtain \eqref{ebi-top}.
\end{proof}

This proof retains the established Brun--Titchmarsh input but uses no
logarithmic-form estimate. Its exceptional roots and the signed
prime costs remain uncontrolled. Full mass at an exponent-one prime
is positive while its signed cost is $-2\log p$; the displayed
concentration therefore asserts neither a radical lower bound nor a
failure of the signed-tail gate. The ordinary proofs were independently
reviewed. A separate exact replay checks 571 actual roots and 1142
valuation equalities, without using floating-point angles as proof.
The present estimates are not claimed as Lean formalizations.
